\documentclass[../../../include/open-logic-section]{subfiles}

\begin{document}

\olfileid{sth}{cardinals}{zfc}
\olsection{$\ZFC$: معلم رئيس}

بإضافة الترتيب الجيد نكون قد بلغنا آخر المعالم النظرية الرئيسة. فقد
أصبحت لدينا الآن جميع البديهيات اللازمة لنظرية~$\ZFC$. وتفصيل ذلك كما
يأتي:

\begin{defn}
لنظرية $\ZFC$ البديهيات الآتية: الامتدادية، والاتحاد، والأزواج، ومجموعات
القوى، واللانهاية، والأساس، والترتيب الجيد، وجميع حالات مخططي الفصل
والاستبدال. وبعبارة أخرى، تضيف $\ZFC$ الترتيب الجيد إلى~$\ZF$.
\end{defn}

يرمز $\ZFC$ إلى نظرية المجموعات \emph{لزيرميلو--فرينكل} مع
\emph{الاختيار}. وقد يبدو هذا غريباً بعض الشيء، لأن البديهية التي أضفناها
تسمى «الترتيب الجيد» لا «الاختيار». لكن حين نصوغ {الاختيار} لاحقاً سيتبين
أن الترتيب الجيد يكافئ الاختيار (بالنسبة إلى~$\ZF$؛ انظر
\olref[choice][woproblem]{thmwochoice}). ومن ثم لا فرق في أيهما نتخذه
بديهية «أساسية». كما أن الاسم «$\ZFC$» هو الاسم القياسي تماماً في
الأدبيات.

\end{document}
